% Compile with xelatex (twice).
% Talk script v3 (2026-10-18) — final specification: 2 models (binary M1,
% multinomial M2), 8 unique variables, no ABSA, no keyword count, no reactions.
% Aligned 1-to-1 with presentation_v2.pptx (21 slides). Target ~15 minutes.
%
\documentclass[10pt,a4paper]{article}

\usepackage{fontspec}
\usepackage{xeCJK}
\setCJKmainfont[AutoFakeSlant=0.15]{Malgun Gothic}
\setCJKsansfont[AutoFakeSlant=0.15]{Malgun Gothic}
\xeCJKsetup{CJKspace=true}

\usepackage[top=1.4cm, bottom=1.4cm, left=1.4cm, right=1.4cm]{geometry}
\usepackage{booktabs}
\usepackage{url}
\usepackage{enumitem}
\usepackage{microtype}
\usepackage{amsmath}

\linespread{1.1}
\setlength{\parindent}{0pt}
\setlength{\parskip}{3pt}

% ── layout helpers ────────────────────────────────────────────────────────────
\newcommand{\sechead}[2]{%
  \vspace{6pt}%
  \noindent\rule{\linewidth}{0.6pt}\\[1pt]
  {\small\bfseries #1\hfill #2}\\[-3pt]
  \rule{\linewidth}{0.4pt}\par\vspace{2pt}%
}

\newcommand{\sent}[1]{#1\\[5pt]}

\newcommand{\bilingual}[2]{%
  \noindent%
  \begin{minipage}[t]{0.485\linewidth}
    \small #1
  \end{minipage}\hfill\vrule\hfill%
  \begin{minipage}[t]{0.483\linewidth}
    \small #2
  \end{minipage}\par\vspace{4pt}%
}

% ── document ─────────────────────────────────────────────────────────────────
\begin{document}

\begin{center}
  {\large\bfseries Talk Script --- Workshop paper}\\[3pt]
  {\large\bfseries 발표 스크립트 --- 워크숍 논문}\\[2pt]
  {\small \textit{An Empirical Analysis of Visit Context and Review Outcomes
  in Online Reviews of Korean Tourist Attractions}}\\[2pt]
  {\small Aligned with \texttt{presentation\_v2.pptx} (21 slides) ---
  target $\sim$15 minutes}
\end{center}

\vspace{4pt}

% ═════════════════════════════════════════════════════════════════════════════
\sechead{Slide 1 --- Title}{$\sim$30s}

\bilingual{%
\sent{Good morning, everyone. Thank you for being here.}
\sent{My name is Jorge Soriano, from the Department of Data Science
at Hoseo University.}
\sent{Today I will present an empirical analysis of how visit context
shapes online reviews of Korean tourist attractions on Naver Place.}
}{%
\sent{안녕하세요, 여러분. 참석해 주셔서 감사합니다.}
\sent{저는 호서대학교 데이터사이언스학과의 호르헤 소리아노입니다.}
\sent{오늘 저는 방문 맥락이 한국 관광지에 대한 네이버 플레이스 온라인
리뷰에 어떻게 영향을 주는지에 대한 실증 분석을 발표하겠습니다.}
}

% ═════════════════════════════════════════════════════════════════════════════
\sechead{Slide 2 --- Agenda}{$\sim$30s}

\bilingual{%
\sent{The talk has five parts.}
\sent{First, I introduce the research background and the two research
questions.}
\sent{Second, I describe the data and the two statistical models.}
\sent{Third, I present the results for each model.}
\sent{Fourth, I interpret the two models jointly.}
\sent{Finally, I close with implications and limitations.}
}{%
\sent{발표는 다섯 부분으로 구성됩니다.}
\sent{첫째, 연구 배경과 두 가지 연구 질문을 소개합니다.}
\sent{둘째, 데이터와 두 가지 통계 모형을 설명합니다.}
\sent{셋째, 각 모형의 결과를 제시합니다.}
\sent{넷째, 두 모형을 함께 해석합니다.}
\sent{마지막으로, 시사점과 한계를 다룹니다.}
}

% ═════════════════════════════════════════════════════════════════════════════
\sechead{Slide 3 --- Section Title: Introduction}{$\sim$10s}

\bilingual{%
\sent{Let's start with the introduction.}
}{%
\sent{먼저 서론부터 시작하겠습니다.}
}

% ═════════════════════════════════════════════════════════════════════════════
\sechead{Slide 4 --- Research background}{$\sim$60s}

\bilingual{%
\sent{Travellers routinely consult online reviews before deciding where
to visit or whether to return.}
\sent{Most review platforms expose only the review text and the star
rating.}
\sent{But Naver Place is different: it also records structured metadata
about the visit itself --- who you went with, why you went, whether you
had to wait at the entrance, and whether you made a reservation.}
\sent{Prior work either analyses review text in isolation or aggregates
reviews up to the business level, losing the review as the unit of
analysis.}
\sent{Lee, Jeong and Kim's 2026 study, which uses Naver Smart Place at
the business--week level, is a key reference, but it does not observe
individual visit context.}
\sent{Our paper fills that gap.}
}{%
\sent{여행자는 어디를 방문하거나 재방문할지 결정하기 전에 일상적으로
온라인 리뷰를 참고합니다.}
\sent{대부분의 리뷰 플랫폼은 리뷰 텍스트와 별점만 공개합니다.}
\sent{하지만 네이버 플레이스는 다릅니다. 누구와 갔는지, 왜 갔는지,
입구에서 대기했는지, 예약했는지에 대한 구조화된 메타데이터도 기록합니다.}
\sent{기존 연구는 리뷰 텍스트만 분석하거나 리뷰를 업체 단위로 집계해
리뷰 자체를 분석 단위로 다루지 못합니다.}
\sent{이민영·정동영·김지영 교수의 2026년 연구는 중요한 참조이지만, 업체-주
단위 자료를 쓰기 때문에 개별 방문 맥락은 관측되지 않습니다.}
\sent{본 연구는 바로 그 공백을 메웁니다.}
}

% ═════════════════════════════════════════════════════════════════════════════
\sechead{Slide 5 --- Research questions}{$\sim$60s}

\bilingual{%
\sent{We ask two research questions.}
\sent{RQ1 is about Revisit: how do visit context and reviewer
characteristics shape the probability that a review is written by a
returning visitor?}
\sent{RQ2 is about Rating: how do the same factors shape the probability
of each star-rating category --- top with five stars, mid with four to
four-and-a-half, and low with three-and-a-half or fewer?}
\sent{Our framework treats both outcomes as observable traces of a
latent satisfaction construct.}
\sent{The rating is a declared trace; the revisit flag is a behavioural
trace.}
\sent{We model them in parallel with mostly the same predictors, so that
differences in coefficients reflect differences in the outcomes
themselves, not in how we estimate them.}
}{%
\sent{우리는 두 가지 연구 질문을 다룹니다.}
\sent{첫 번째는 재방문에 관한 것입니다. 방문 맥락과 작성자 특성이
리뷰어가 재방문자일 확률에 어떻게 영향을 미치는가입니다.}
\sent{두 번째는 별점에 관한 것입니다. 같은 요인들이 세 범주의 별점
---5점의 TOP, 4~4.5점의 MID, 3.5점 이하의 LOW--- 확률을 어떻게 결정하는가입니다.}
\sent{우리의 틀은 두 결과를 모두 잠재적 만족도 개념의 관측 가능한
흔적으로 간주합니다.}
\sent{별점은 선언된 흔적이고, 재방문 표식은 행동적 흔적입니다.}
\sent{거의 동일한 예측 변수로 두 모형을 병렬 추정합니다. 그래야 계수의
차이가 모형 설정의 차이가 아니라 결과 변수 자체의 차이에서 나온다고
해석할 수 있습니다.}
}

% ═════════════════════════════════════════════════════════════════════════════
\sechead{Slide 6 --- Section Title: Methodology}{$\sim$10s}

\bilingual{%
\sent{Let's move to the methodology.}
}{%
\sent{이제 연구 방법으로 넘어가겠습니다.}
}

% ═════════════════════════════════════════════════════════════════════════════
\sechead{Slide 7 --- Data}{$\sim$45s}

\bilingual{%
\sent{The dataset consists of 14{,}401 Korean-language reviews from
8{,}276 reviewers across 29 attractions.}
\sent{Reviews were collected by scraping Naver Place for each attraction
individually.}
\sent{Of the 14{,}401 reviews, 11{,}165 carry free text; these form the
sample for the revisit model.}
\sent{Among those, 4{,}044 also carry a star rating; these form the
sample for the rating model.}
\sent{The 29 attractions were selected by stratified sampling across
typological categories --- palaces, traditional quarters, markets,
museums, parks, beaches, theme parks --- with a minimum of thirty
Korean reviews per site.}
}{%
\sent{자료는 29개 관광지의 8{,}276명 작성자가 작성한 14{,}401건의 한국어
리뷰로 구성됩니다.}
\sent{리뷰는 네이버 플레이스에서 각 관광지별로 스크래핑하여 수집했습니다.}
\sent{14{,}401건 중 11{,}165건은 리뷰 텍스트를 포함하며, 재방문 모형의
표본이 됩니다.}
\sent{그 중 4{,}044건은 별점도 함께 가지며, 별점 모형의 표본이 됩니다.}
\sent{29개 관광지는 궁궐, 전통마을, 시장, 박물관, 공원, 해변, 테마파크
등 유형별 범주에 걸쳐 층화 추출로 선정했으며, 각 관광지는 최소 30건
이상의 한국어 리뷰를 포함합니다.}
}

% ═════════════════════════════════════════════════════════════════════════════
\sechead{Slide 8 --- Variables table}{$\sim$60s}

\bilingual{%
\sent{The table on the slide lists our predictor variables.}
\sent{Three variables are shared across both models: sentiment of the
text, waiting at the entrance, and the log of reviewer followers.}
\sent{These three are the only predictors that are both theoretically
relevant and statistically significant in both outcomes.}
\sent{Model 1 adds three predictors that matter for revisit but not for
rating: tourist indicator, solo visit indicator, and reviewer's prior
review count.}
\sent{Model 2 adds two predictors that matter for rating but not for
revisit: number of photos attached to the review, and a dummy for
autumn visits.}
\sent{In total we use eight unique variables across the two models.}
}{%
\sent{슬라이드의 표는 예측 변수를 보여줍니다.}
\sent{세 변수는 두 모형 모두에 공유됩니다. 텍스트의 감성, 입구에서의
대기 여부, 그리고 리뷰어 팔로워 수의 로그값입니다.}
\sent{이 세 변수만이 이론적으로 관련이 있고 두 결과 모두에 통계적으로
유의한 유일한 예측 변수입니다.}
\sent{모형 1은 재방문에는 유의하지만 별점에는 그렇지 않은 세 변수를
추가합니다. 관광 목적 지시 변수, 혼자 방문 지시 변수, 그리고 작성자의
과거 리뷰 수입니다.}
\sent{모형 2는 별점에는 유의하지만 재방문에는 그렇지 않은 두 변수를
추가합니다. 리뷰에 첨부된 사진 수와 가을 방문 더미입니다.}
\sent{두 모형에 걸쳐 총 여덟 개의 고유 변수를 사용합니다.}
}

% ═════════════════════════════════════════════════════════════════════════════
\sechead{Slide 9 --- Model 1 equation}{$\sim$45s}

\bilingual{%
\sent{Model 1 is a standard binary logistic regression.}
\sent{The dependent variable is revisit, equal to one if the reviewer
has visited at least twice.}
\sent{The right-hand side includes six predictors --- the three shared
variables plus three M1-specific ones --- and a place fixed effect delta
sub j for each of the 29 attractions.}
\sent{Standard errors are clustered by reviewer, since the same person
often writes several reviews.}
}{%
\sent{모형 1은 표준적인 이항 로지스틱 회귀입니다.}
\sent{종속 변수는 재방문이며, 리뷰어가 최소 두 번 이상 방문한 경우
1입니다.}
\sent{우변에는 여섯 개의 예측 변수---세 개의 공유 변수와 세 개의 M1
전용 변수---와 29개 관광지 각각에 대한 장소 고정효과 델타 서브 j가
들어갑니다.}
\sent{표준오차는 같은 사람이 여러 리뷰를 쓰는 경우가 많기 때문에 리뷰어
수준에서 군집 표준오차를 사용합니다.}
}

% ═════════════════════════════════════════════════════════════════════════════
\sechead{Slide 10 --- Model 2 equation}{$\sim$60s}

\bilingual{%
\sent{Model 2 is a multinomial logistic regression on the three-category
rating.}
\sent{The reference category is TOP --- five stars.}
\sent{For each non-reference category k, that is MID and LOW, we model
the log-odds of being in that category relative to TOP as a function of
five predictors --- the three shared variables plus two M2-specific ones.}
\sent{The multinomial specification gives us two parallel sets of
coefficients: one set for MID versus TOP, another for LOW versus TOP.}
\sent{Each category has its own intercept, its own slope vector, and its
own place fixed effects.}
}{%
\sent{모형 2는 3범주 별점에 대한 다항 로지스틱 회귀입니다.}
\sent{참조 범주는 TOP, 즉 5점입니다.}
\sent{각 비참조 범주 k, 즉 MID와 LOW에 대해, 해당 범주가 TOP 대비일
로그 오즈를 다섯 개 예측 변수---세 개의 공유 변수와 두 개의 M2 전용
변수---의 함수로 모형화합니다.}
\sent{다항 설정은 두 세트의 평행한 계수를 제공합니다. 하나는 MID 대
TOP, 다른 하나는 LOW 대 TOP입니다.}
\sent{각 범주는 고유의 절편, 고유의 기울기 벡터, 고유의 장소 고정효과를
가집니다.}
}

% ═════════════════════════════════════════════════════════════════════════════
\sechead{Slide 11 --- Why multinomial?}{$\sim$60s}

\bilingual{%
\sent{Why multinomial and not ordered logit?}
\sent{Initially we tried an ordered logistic regression, which is the
natural choice for an ordinal outcome like a star rating.}
\sent{But the Brant test rejected the proportional-odds assumption:
chi-squared is 18.8, 7 degrees of freedom, p-value 0.009.}
\sent{Two predictors violated the assumption individually: sentiment and
the log of review age.}
\sent{Multinomial logit removes this assumption by design: each category
comparison gets its own coefficients, with no equality constraint.}
\sent{The trade-off is that the ordinal structure of the rating is not
imposed on the fit, but our within-plus-or-minus-one accuracy of 95
percent confirms the ordering is still captured empirically.}
}{%
\sent{왜 순서형 로지스틱이 아니라 다항 로지스틱인가요?}
\sent{처음에는 순서형 로지스틱을 시도했습니다. 별점 같은 순서형 결과에
자연스러운 선택이기 때문입니다.}
\sent{하지만 Brant 검정에서 비례오즈 가정이 기각되었습니다. 카이제곱
18.8, 자유도 7, p값 0.009입니다.}
\sent{두 변수가 개별적으로 가정을 위반했습니다. 감성과 리뷰 연령의
로그값입니다.}
\sent{다항 로지스틱은 설계상 이 가정을 제거합니다. 각 범주 비교가
동일성 제약 없이 자체 계수를 가집니다.}
\sent{트레이드오프는 별점의 순서적 구조가 적합에 강제되지 않는다는
것이지만, 95\%의 within-$\pm 1$ 정확도는 순서가 여전히 경험적으로
포착되고 있음을 확인시켜 줍니다.}
}

% ═════════════════════════════════════════════════════════════════════════════
\sechead{Slide 12 --- Section Title: Results}{$\sim$10s}

\bilingual{%
\sent{Let's look at the results.}
}{%
\sent{이제 결과를 보겠습니다.}
}

% ═════════════════════════════════════════════════════════════════════════════
\sechead{Slide 13 --- Model fit summary}{$\sim$45s}

\bilingual{%
\sent{This table summarises the fit of both models.}
\sent{Model 1 has pseudo-R-squared of 0.20, AIC of 9{,}602.}
\sent{Model 2 has pseudo-R-squared of 0.19, AIC of 4{,}754.}
\sent{Both are in the "good" range for discrete-choice models in
consumer research.}
\sent{Out-of-sample performance from five-fold cross-validation grouped
by reviewer: Model 1 AUC is 0.80, Model 2 exact accuracy is 0.77 with
95 percent within-one-category accuracy.}
}{%
\sent{이 표는 두 모형의 적합도를 요약합니다.}
\sent{모형 1은 pseudo-R-squared 0.20, AIC 9{,}602입니다.}
\sent{모형 2는 pseudo-R-squared 0.19, AIC 4{,}754입니다.}
\sent{두 모형 모두 소비자 연구의 이산 선택 모형에서 ``좋음'' 범위에
속합니다.}
\sent{리뷰어 기준 5-fold 교차검증 결과, 모형 1의 AUC는 0.80, 모형 2의
정확 정확도는 0.77이며 within-±1 정확도는 95\%입니다.}
}

% ═════════════════════════════════════════════════════════════════════════════
\sechead{Slide 14 --- Model 1 Results (Revisit)}{$\sim$90s}

\bilingual{%
\sent{Here are the Model 1 coefficients and the forest plot.}
\sent{Six predictors, all significant at p less than 0.01.}
\sent{The strongest negative effect is waiting at the entrance ---
odds ratio 0.54.}
\sent{In plain terms, a reviewer who waited has about half the odds of
being a return visitor compared to one who did not wait.}
\sent{This is our clearest managerial finding: waiting is the single
operational variable in this model that attraction operators can
directly reduce.}
\sent{Two reviewer-identity variables enter positively: log of prior
reviews --- OR 1.31 --- and log of followers --- OR 1.23.}
\sent{More active or more-connected reviewers are disproportionately
likely to be returning visitors; this is probably self-selection, but
it is informative.}
\sent{Visit-type effects are large and in opposite directions.}
\sent{Tourists are 45 percent less likely to return --- OR 0.55 ---
consistent with one-time visits.}
\sent{Solo visitors are more than twice as likely to return --- OR 2.26
--- probably locals or committed enthusiasts.}
\sent{Sentiment is modestly negative, OR 0.74.}
\sent{One interpretation: first-time visitors write uniformly positive
reviews, while repeat visitors introduce more nuance.}
}{%
\sent{여기 모형 1의 계수와 forest plot이 있습니다.}
\sent{여섯 개의 예측 변수 모두 p 0.01 수준에서 유의합니다.}
\sent{가장 강한 부정적 효과는 입구 대기---오즈비 0.54---입니다.}
\sent{쉽게 말해, 대기한 리뷰어는 대기하지 않은 사람에 비해 재방문자일
확률이 약 절반입니다.}
\sent{이것이 가장 명확한 운영적 발견입니다. 대기는 관광지 운영자가 직접
줄일 수 있는 유일한 운영 변수입니다.}
\sent{두 작성자 정체성 변수는 양의 방향입니다. 과거 리뷰 수의 로그값
---OR 1.31---과 팔로워 수의 로그값---OR 1.23---입니다.}
\sent{더 활발하거나 더 연결된 리뷰어는 재방문자일 가능성이 더 높습니다.
자기선택일 가능성이 크지만 의미 있는 정보입니다.}
\sent{방문 유형 효과는 크고 반대 방향입니다.}
\sent{관광객은 재방문할 확률이 45\% 낮습니다---OR 0.55---. 일회성
방문과 일치합니다.}
\sent{혼자 방문자는 재방문할 확률이 두 배 이상입니다---OR 2.26---.
지역 주민이거나 열성 애호가일 가능성이 큽니다.}
\sent{감성은 약하게 음의 방향입니다, OR 0.74.}
\sent{한 가지 해석은 신규 방문자가 균일하게 긍정적인 리뷰를 쓰는 반면,
재방문자는 더 세밀한 평가를 담는다는 것입니다.}
}

% ═════════════════════════════════════════════════════════════════════════════
\sechead{Slide 15 --- Model 2 Results (Rating)}{$\sim$90s}

\bilingual{%
\sent{Now Model 2 --- the rating multinomial.}
\sent{The table has two columns: MID vs TOP on the left, LOW vs TOP on
the right.}
\sent{The key insight is that predictors operate asymmetrically across
these two comparisons.}
\sent{Sentiment is dominant in both --- OR 0.44 for MID, OR 0.04 for
LOW --- but much stronger for LOW.}
\sent{A note of caution: because sentiment and rating are produced in
the same reviewing session, this is partly mechanical, not causal.}
\sent{Waiting shows a striking asymmetric pattern: not significant for
MID vs TOP, but it TRIPLES the odds of LOW vs TOP --- OR 3.10.}
\sent{In other words, waiting is a disaster signal, not a mediocrity
signal.}
\sent{Reviewer influence --- log of followers --- is only significant
for LOW vs TOP, with OR 0.71: more-followed reviewers are protected
against very low ratings.}
\sent{Season autumn has a strong, symmetric effect: OR about 0.4 for
both MID and LOW --- autumn visits are strongly concentrated in TOP
ratings, consistent with autumn being Korea's peak tourism season for
foliage.}
}{%
\sent{이제 모형 2---별점 다항 로지스틱입니다.}
\sent{표에는 두 열이 있습니다. 왼쪽은 MID 대 TOP, 오른쪽은 LOW 대
TOP입니다.}
\sent{핵심 통찰은 예측 변수가 두 비교에서 비대칭적으로 작용한다는
것입니다.}
\sent{감성은 두 비교 모두에서 지배적입니다---MID에서 OR 0.44, LOW에서
OR 0.04---하지만 LOW에서 훨씬 강합니다.}
\sent{주의할 점은 감성과 별점이 동일한 리뷰 세션에서 생성되므로, 이는
부분적으로 기계적이며 인과적이지 않다는 것입니다.}
\sent{대기는 두드러진 비대칭 패턴을 보입니다. MID 대 TOP에는 유의하지
않지만, LOW 대 TOP의 오즈를 세 배로 늘립니다---OR 3.10.}
\sent{다시 말해, 대기는 ``실패 신호''이지 ``평범함 신호''가 아닙니다.}
\sent{리뷰어 영향력---팔로워 수의 로그값---은 LOW 대 TOP에만 유의하며
OR 0.71입니다. 팔로워가 많은 리뷰어는 매우 낮은 별점으로부터
보호받습니다.}
\sent{가을 효과는 강하고 대칭적입니다. MID와 LOW 모두에서 OR 약
0.4입니다. 가을 방문은 TOP 별점에 강하게 집중되어 있으며, 가을이 한국
단풍 관광의 성수기임과 일치합니다.}
}

% ═════════════════════════════════════════════════════════════════════════════
\sechead{Slide 16 --- Forest plot M2}{$\sim$45s}

\bilingual{%
\sent{This visualisation makes the asymmetry clear.}
\sent{Each row is one predictor. The left panel is MID vs TOP, the
right panel is LOW vs TOP. Blue markers are significant at p less than
0.05.}
\sent{Notice how waited barely moves in the left panel but swings
strongly in the right.}
\sent{Sentiment also swings farther in the LOW panel.}
\sent{This is exactly the pattern the ordered logit could not accommodate,
which is why we chose the multinomial specification.}
}{%
\sent{이 시각화는 비대칭성을 명확하게 보여줍니다.}
\sent{각 행은 하나의 예측 변수입니다. 왼쪽 패널은 MID 대 TOP, 오른쪽은
LOW 대 TOP입니다. 파란색 마커는 p 0.05 수준에서 유의한 변수입니다.}
\sent{왼쪽 패널에서는 거의 움직이지 않는 ``대기''가 오른쪽에서는 강하게
움직인다는 점을 주목해 주십시오.}
\sent{감성도 LOW 패널에서 훨씬 크게 움직입니다.}
\sent{이것이 바로 순서형 로지스틱이 담지 못하는 패턴이며, 다항 설정을
선택한 이유입니다.}
}

% ═════════════════════════════════════════════════════════════════════════════
\sechead{Slide 17 --- Divergent traces of satisfaction}{$\sim$90s}

\bilingual{%
\sent{Now the key interpretation.}
\sent{Revisit and Rating correlate only weakly in our data: r equals
0.08.}
\sent{In other words, the two outcomes share only about zero-point-six
percent of their variance.}
\sent{This matters because, if the two outcomes were essentially the
same construct, we would be running the same regression twice.}
\sent{But they are not. The shared predictors operate differently in
each model.}
\sent{Waited reduces revisit by half, and triples the odds of a bad
rating --- same direction, different magnitude.}
\sent{Sentiment weakly reduces revisit but strongly protects against low
ratings --- different roles entirely.}
\sent{Follower count predicts revisit positively, and protects against
low ratings --- convergent effect.}
\sent{The takeaway: declared satisfaction, the rating, and behavioural
satisfaction, the revisit, are related but genuinely distinct traces of
a latent satisfaction construct.}
\sent{Our two parallel models reveal that different predictors route
through to the two traces differently --- that is the empirical
contribution of this paper.}
}{%
\sent{이제 핵심 해석입니다.}
\sent{우리 자료에서 재방문과 별점은 약한 상관만을 보입니다. r은
0.08입니다.}
\sent{다시 말해, 두 결과는 분산의 약 0.6\%만을 공유합니다.}
\sent{이 점이 중요한 이유는, 만약 두 결과가 본질적으로 같은 구성개념이라면,
우리는 같은 회귀를 두 번 돌리는 셈이기 때문입니다.}
\sent{하지만 그렇지 않습니다. 공유 예측 변수가 각 모형에서 다르게 작용합니다.}
\sent{대기는 재방문을 절반으로 줄이고, 낮은 별점의 오즈를 세 배로
늘립니다---방향은 같지만 크기가 다릅니다.}
\sent{감성은 재방문을 약하게 줄이지만 낮은 별점을 강하게 방지합니다
---역할이 완전히 다릅니다.}
\sent{팔로워 수는 재방문을 양의 방향으로 예측하고 낮은 별점을 방지합니다
---수렴 효과입니다.}
\sent{요점은, 선언된 만족도인 별점과 행동적 만족도인 재방문이, 관련은
있지만 잠재적 만족도 개념의 진정으로 구별되는 흔적이라는 점입니다.}
\sent{우리의 두 병렬 모형은 서로 다른 예측 변수가 두 흔적으로 다르게
경로화됨을 보여줍니다---이것이 본 논문의 경험적 기여입니다.}
}

% ═════════════════════════════════════════════════════════════════════════════
\sechead{Slide 18 --- Section Title: Conclusion}{$\sim$10s}

\bilingual{%
\sent{Finally, let me close with implications and limitations.}
}{%
\sent{마지막으로 시사점과 한계로 마무리하겠습니다.}
}

% ═════════════════════════════════════════════════════════════════════════════
\sechead{Slide 19 --- Summary and implications}{$\sim$60s}

\bilingual{%
\sent{Four main findings.}
\sent{First, waiting at the entrance is the strongest and most
consistent negative predictor: it halves revisit odds and triples odds
of a LOW rating.}
\sent{Second, reviewer-identity variables predict revisit well; weakly
for rating.}
\sent{Third, tourists and solo visitors enter revisit with opposite
signs: tourists down, alone visitors up.}
\sent{Fourth, autumn visits concentrate heavily in the TOP rating
category.}
\sent{For attraction managers, the actionable implication is: reduce
entrance wait times --- it is the one lever that improves both revisit
and rating.}
\sent{For platforms, investment in reviewer retention likely outperforms
content-side nudges for driving repeat visits.}
}{%
\sent{네 가지 주요 발견입니다.}
\sent{첫째, 입구 대기는 가장 강하고 일관된 부정적 예측 변수입니다.
재방문 오즈를 절반으로 줄이고 LOW 별점 오즈를 세 배로 늘립니다.}
\sent{둘째, 작성자 정체성 변수는 재방문을 잘 예측합니다. 별점에는
약합니다.}
\sent{셋째, 관광객과 혼자 방문자는 재방문에서 반대 방향입니다. 관광객은
감소, 혼자 방문자는 증가입니다.}
\sent{넷째, 가을 방문은 TOP 별점 범주에 집중됩니다.}
\sent{관광지 운영자에게 실행 가능한 시사점은, 입구 대기 시간을 줄이라는
것입니다. 재방문과 별점 모두를 개선하는 유일한 레버입니다.}
\sent{플랫폼의 경우, 리뷰어 유지에 대한 투자가 재방문을 유도하는
콘텐츠 측 넛지보다 더 효과적일 가능성이 큽니다.}
}

% ═════════════════════════════════════════════════════════════════════════════
\sechead{Slide 20 --- Limitations and future work}{$\sim$45s}

\bilingual{%
\sent{Five limitations to be clear about.}
\sent{One: associations only. The design is cross-sectional, we cannot
identify causal effects.}
\sent{Two: Model 1 profiles returning visitors; it does not explain why
anyone returns, because the revisit label is set before the review is
written.}
\sent{Three: the sentiment-rating link is partly mechanical, since both
come from the same reviewing session.}
\sent{Four: the MID rating class is small and sandwiched between TOP and
LOW; the multinomial routes most predictions to the extremes.}
\sent{Five: single platform --- we cannot separate Naver's reviewer
population from Korean reviewers in general.}
\sent{For future work: add month fixed effects, run a Hausman test, and
seek instrumental variables for waiting.}
}{%
\sent{분명히 밝힐 다섯 가지 한계입니다.}
\sent{첫째, 연관 분석에 한정됩니다. 횡단면적 설계이므로 인과적 효과를
식별할 수 없습니다.}
\sent{둘째, 모형 1은 재방문자를 프로파일링합니다. 재방문 표식은 리뷰
작성 이전에 결정되므로, 왜 재방문하는지는 설명하지 못합니다.}
\sent{셋째, 감성-별점 연관은 두 변수 모두 같은 리뷰 세션에서 나오므로
부분적으로 기계적입니다.}
\sent{넷째, MID 별점 범주는 작고 TOP과 LOW 사이에 끼어 있어, 다항
모형은 대부분의 예측을 극단으로 라우팅합니다.}
\sent{다섯째, 단일 플랫폼입니다. 네이버의 리뷰어 모집단을 일반 한국
리뷰어와 분리할 수 없습니다.}
\sent{향후 과제로는 월 고정효과 추가, Hausman 검정 수행, 그리고 대기에
대한 도구 변수 탐색이 있습니다.}
}

% ═════════════════════════════════════════════════════════════════════════════
\sechead{Slide 21 --- Thank you}{$\sim$15s}

\bilingual{%
\sent{Thank you very much for your attention.}
\sent{I welcome your questions and comments.}
}{%
\sent{경청해 주셔서 대단히 감사합니다.}
\sent{질문과 코멘트를 환영합니다.}
}

\vspace{10pt}

\noindent\rule{\linewidth}{0.5pt}

\vspace{6pt}

\begin{center}
\small
\textbf{Total estimated time:} $\sim$15 minutes (based on $\sim$120--140 words
per minute for a presentation delivered at a comfortable pace).\\[2pt]
\textbf{총 예상 시간:} 약 15분 (분당 120-140 단어 기준).
\end{center}

\end{document}
